\section{Dérivation d'ordre supérieurs}

\subsection{Dérivées secondes}
\newdef{Soit $E \subset \R$ ouvert non-vide et $f : E \to \R$ telle que toutes les dérivées partielles secondes $\frac{\partial^2 f}{\partial x_i \partial x_j} : E \to \R, ~i,j \in \entiers$ existent. On appelle matrice hessienne de $f$ en $\bm x \in E$ la matrice $H_f(\bm x) \in \R^{n \times n}$ 
	\[
	H_f(\bm x) = \begin{bmatrix}
		\frac{\partial^2 f}{\partial x_1^2}(\bm x) & \frac{\partial^2 f}{\partial x_1 \partial x_2} (\bm x) & \ldots & \frac{\partial^2 f}{\partial x_1 \partial x_2} (\bm x) \\
		\vdots & & & \vdots \\
		\frac{\partial^2 f}{\partial x_n \partial x_1} (\bm x) & \ldots & \ldots & \frac{\partial^2 f}{\partial x_n^2}(\bm x)
	\end{bmatrix}
	\]
	De manière plus concise on a,
	\[
	\left(H_f(\bm x)\right)_{ij} = \frac{\partial^2 f}{\partial x_i \partial x_j} (\bm x)
	\]
	Si toutes dérivées partielles secondes $ \frac{\partial^2 f}{\partial x_i \partial x_j} (\bm x) : E \to \R$ existent et sont continues sur $E$, alors $f \in C^2$.
}{}

\newlem{Soit $f \in C^2(E)$ avec $E \subset \rn$ un ouvert non-vide et $\bm v, \bm w \in \rn$ unitaires. Alors
	\[
	D_{\bm w \bm v}^2 f(\bm x) = \bm w^\top H_f(\bm x) \bm v,
	\]
}{}

\prf{Puisque $f \in C^2(E)$, alors $D_{\bm v} f \in C^1(E)$ et donc
\[
D_{\bm w} (D_{\bm v}f)(\bm x) ) \sumn \frac{\partial (D_{\bm v} f)}{\partial x_i}(\bm x) w_i = \sum_{i,j = 1}^n \frac{\partial^2 f}{\partial x_i \partial x_j}(\bm x) v_j w_i = \bm w^\top H_f(\bm x) \bm v
\]
}
 

\newthm{Soit $f : E \to \R$ telle que $\frac{\partial f}{\partial x_i}, \frac{\partial f}{\partial x_j}, \frac{\partial^2 f}{\partial x_i \partial x_j}, \frac{\partial^2 f}{\partial x_j \partial x_i}$ existent dans $E$ et $\frac{\partial^2 f}{\partial x_i \partial x_j}, \frac{\partial^2 f}{\partial x_j \partial x_i}$ sont continues en $\bm x \in E$, alors
	\[
	\frac{\partial^2 f}{\partial x_i \partial x_j}(\bm x) = \frac{\partial^2 f}{\partial x_j \partial x_i}(\bm x)
	\]
}{Théorème de Schwarz}

\prf{On considère la quantité,
	\[
	\Delta(s,t) = f(\bm x + s \bm e_i + t \bm e_j) - f(\bm x + \bm e_i) - f(\bm x + t \bm e_j) + f(\bm x)
	\]
et les deux fonctions auxiliaires,
\[
g(\xi) = f(\bm x + \xi \bm e_i + t \bm e_j) - f(\bm x + \xi \bm e_i) \et h(\xi) = f(\bm x + s \bm e_i + \xi \bm e_j) - f(\bm x + \xi \bm e_j)
\]
La fonction $g$ est dérivable sur $]0,s[$ et donc par le théorème des accroissement finis, il existe $\tilde s \in ]0,s[$ tel que
\[
g(s) - g(0) = g'(\tilde s) s = \left( \frac{\partial f}{\partial x_i}(\bm x + \tilde s \bm e_i + t \bm e_j) - \frac{\partial f}{\partial x_i}(\bm x + \tilde s \bm e_i) \right) s.
\]
Si on pose
\[
\varphi(\eta) = \frac{\partial f}{\partial x_i}(\bm x + \tilde s \bm e_i + \eta \bm e_j)
\]
la fonction $\varphi$ est dérivable sur $]0,t[$ donc il existe $\tilde t \in ]0,t[$ tel que
\[
\varphi(t) - \varphi(0) = \varphi'(\tilde t)t = \frac{\partial^2 f}{\partial x_j \partial x_i} (\bm x + \tilde s \bm e_i + \tilde t \bm e_j) t
\]
On a finalement,
\[
\Delta(s,t) = g(s) - g(0) = g'(\tilde s) s = \left( \varphi(t) - \varphi(0)\right) s = \varphi(\tilde t) st = \frac{\partial^2 f}{\partial x_j \partial x_i} (\bm x + \tilde s \bm e_i + \tilde t \bm e_j) st
\]
De manière analogue avec la fonction $h$ on trouve
\[
\Delta(s,t) = \frac{\partial^2 f}{\partial x_i \partial x_j} (\bm x + \hat s \bm e_i + \hat t \bm e_j) st
\]
et donc 
\[
\Delta(s,t) = \frac{\partial^2 f}{\partial x_i \partial x_j} (\bm x + \hat s \bm e_i + \hat t \bm e_j) st = \varphi(\tilde t) st = \frac{\partial^2 f}{\partial x_j \partial x_i} (\bm x + \tilde s \bm e_i + \tilde t \bm e_j) st
\]
Si on prend $s = t \to 0^+$, on obtient $\hat s, \tilde s \to 0^+$ et $\hat t, \tilde t \to 0^+$. Et donc par continuité des dérivées partielles secondes en $\bm x$, on a
\[
\lim_{s \to 0^+} \frac{1}{s^2} \Delta(s,s) = \frac{\partial^2 f}{\partial x_i \partial x_j}(\bm x) = \frac{\partial^2 f}{\partial x_j \partial x_i}(\bm x)
\]
}	

\newcor{Soit $E \subset \rn$ ouvert non-vide avec $n \geqslant 2$ et $f \in C^2(E)$. Alors
	\[
	\frac{\partial^2 f}{\partial x_i \partial x_j}(\bm x) = \frac{\partial^2 f}{\partial x_j \partial x_i}(\bm x), \quad \forall \bm x \in E, \quad \forall i,j \in \entiers
	\]
	Le résultat du théorème de Schwarz n'est pas forcément vrai sans l'hypothèse de continuité des dérivées partielles secondes.
}{}

\subsection{Dérivées partielles d'ordres supérieurs à 2}

\newdef{Soit $f : E \to \R$. On généralise la notion de dérivée partielle d'ordre $p$ de $f$ par rapport aux variables $x_{i_1}, x_{i_2}, \ldots, x_{i_p}$ pour $i_k \in \entiers$,
	\[
	\frac{\partial^p f}{\partial x_{i_p} \ldots \partial x_{i_2} \partial x_{i_1}} (\bm x) = \frac{\partial}{\partial x_{i_p}} \left( \dots  \frac{\partial}{\partial x_{i_2}} \left( \frac{\partial f}{\partial x_{i_1}} \right) \right)
	\]
}{}

\newrem{Grâce au théorème de Schwarz, si $f \in C^p(E)$, alors
	\[
	\frac{\partial^p f}{\partial x_{i_p} \dots \partial x_{i_1}} = \frac{\partial^p f}{\partial x_{i_{\sigma(1)}} \dots \partial x_{i_{\sigma(p)}}}
	\]
	où $\sigma $ est une permutation arbitraire. Autrement dit, l'ordre de dérivation n'est pas important si $f \in C^p$ et donc on agencer l'ordre des dérivations tel que
	\[
	\frac{\partial^p f}{\partial x_{i_p} \dots \partial x_{i_1}} = \frac{\partial^p f}{\partial x_1^{\alpha_1} \ldots \partial x_n^{\alpha_n}}
	\]
	avec $\alpha_i \in \N$ et $\sumn \alpha_i = p$.
}

\newdef{Soit $\bm \alpha =(\alpha_1, \ldots, \alpha_n) \in \N^n$ tel que $|\bm \alpha | = \sumn \alpha_i = p$ alors on utilise la \emph{notation multi-entiers},
	\[
	\frac{\partial^{|\bm \alpha|} f}{\partial \bm x^{\bm \alpha}} = \frac{\partial^{|\bm \alpha|}f}{\partial x_1^{\alpha_1} \ldots \partial x_n^{\alpha_n}} = \frac{\partial^p f}{\partial x_{i_p} \dots \partial x_{i_1}}
	\]
	où $\partial x_j^{\alpha_j}$ signifie qu'on dérive $\alpha_j$ fois par rapport à $x_j$. Cette notation n'est valable (ou du moins utile) que si $f \in C^p(E)$.
}{}

\subsection{Développement limité et Formule de Taylor}

\newdef{Soit $f : E \to \R$  et $\bm x \in \mathring{E}$. S'il existe un polynôme $q \in \R[y_1, \ldots, y_n]$ tel que 
	\[
	q(\bm y) = \sum_{|\bm \alpha| \leqslant p} c_{\bm \alpha}(\bm y - \bm x)^{\bm \alpha}
	\]
de degré $p$ en $\bm y$ et une fonction $R_p : E \to \R$ tels que
\[
\forall \bm y \in E, \quad f(\bm y) = q(\bm y) + R_p(\bm y)
\]
avec $\lim_{\bm y \to \bm x} \frac{R_p(\bm y)}{\norme{\bm y - \bm x}} = 0$, alors on dit que $f$ admet un développement limité d'ordre $p$ autour de $\bm x$.
}{}

\newthm{Soit $f : E \to \R$ de classe $C^{p+1}$ et $\bm x_0 \in \mathring{E}$ alors $f$ admet le developpement limité d'ordre $p$ en $\bm x_0$ suivant, $\forall \bm x \in E, \quad [\bm x, \bm x_0] \subset E$,
	\[
	f(\bm x) = \sum_{|\bm \alpha| \leqslant p} \frac{1}{\bm \alpha !} \frac{\partial^{|\bm \alpha|} f}{\partial \bm x^{\bm \alpha}}(\bm x) (\bm x - \bm x_0)^{\bm \alpha} + R_p(\bm x)
	\]
	avec les estimations du reste suivants,
	\begin{align*}
	R_p(\bm x) &= \sum_{|\bm \alpha| = p + 1} \frac{1}{\bm \alpha !} \frac{\partial^{p+1}f}{\partial \bm x^{\bm \alpha}} (\bm x + \theta (\bm x - \bm x_0))(\bm x - \bm x_0)^{\bm \alpha} \quad \text{pour un } \theta \in ]0,1[ \\
	&= \sum_{|\bm \alpha| = p + 1} \frac{(p+1)!}{\bm \alpha!} \int_0^1 (1-s)^p \frac{\partial^{p+1}f}{\partial \bm x^{\bm \alpha}} (\bm x + s (\bm x - \bm x_0))(\bm x - \bm x_0)^{\bm \alpha} ~\mathrm ds
	\end{align*}
	
}{}

\prf{Supposons que $[x_0, x] \subset E$, alors on pose
	\[
	\begin{array}{cccc}
		g : & [0,1] & \to & \R \\
			& t & \mapsto & f(\bm x_0 + t(\bm x- \bm x_0))
	\end{array}
	\]
	comme $f$ est supposé dérivable $p$ fois, $g$ l'est aussi. On peut donc utiliser le développement de Taylor pour les fonctions à une seule variables centré en $0$,
	\[
\forall t \in [0,1], \quad 	g(t) = g(0) + g'(0)t + \frac{g''(0)}{2} t^2 + \ldots + \frac{g^{(p)}}{p!}t^p + R_p^g(t)
	\]
	Or, si on note $\bm x_t = \bm x_0 + t(\bm x - \bm x_0)$, 
	\begin{align*}
		g(t) &= f(\bm x_t) \\
		g'(t) &= Df(\bm x_t) (\bm x - \bm x_0) = \sum_{i_1 = 1}^n \frac{\partial f}{\partial x_{i_1}}(\bm x_t)( x_{i_1}-x_{0_{i_1}}) = \sum_{|\bm \alpha| = 1} \frac{\partial^{|\bm \alpha|} f}{\partial \bm x^{\bm \alpha}}(\bm x) (\bm x - \bm x_0)^{\bm \alpha} \\
		g''(t) &= \sum_{i_2 = 1}^n \frac{\partial }{\partial x_{i_2}} \left( \sum_{i_1 = 1}^n \frac{\partial f}{\partial x_{i_1}}(\bm x_t)( x_{i_1}-x_{0_{i_1}}) \right) = \sum_{|\bm \alpha| =2 } \frac{2}{\bm \alpha !} \frac{\partial^{|\bm \alpha|} f}{\partial \bm x^{\bm \alpha}}(\bm x) (\bm x - \bm x_0)^{\bm \alpha}\\
		&\vdots \\
		g^{(p)}(t) &= \sum_{i_1, \ldots, i_p = 1}^n \frac{\partial p}{\partial x_{i_1} \ldots \partial x_{i_p}} (\bm x_t) (x_{i_1}- x_{0_{i_1}}) \dots (x_{i_p}- x_{0_{i_p}}) = \sum_{|\bm \alpha| = p} \frac{p!}{\bm \alpha !} \frac{\partial^{|\bm \alpha|} f}{\partial \bm x^{\bm \alpha}}(\bm x) (\bm x - \bm x_0)^{\bm \alpha}
	\end{align*}
	Finalement, on peut écrire la formule de Taylor
	\[
	f(\bm x) = g(1) = \sum_{k = 0}^p \frac{1}{k!} g^{(k)}(0) + R_p^g(1) = \sum_{|\bm \alpha| \leqslant p} \frac{1}{\bm \alpha !} \frac{\partial^{|\bm \alpha|} f}{\partial \bm x^{\bm \alpha}}(\bm x) (\bm x - \bm x_0)^{\bm \alpha} + R_p(\bm x)
	\]
}

\subsubsection{Formule de Taylor pour des fonctions vectorielles}
Soit $E \subset \rn$ ouvert non-vide et $\bm f \in C^{p+1}(E)$ et $\bm x, \bm x_0 \in E$ tel $[\bm x_0, \bm x] \subset E$. On peut appliquer la formule de Taylor à chaque composante $f_k$ de $\bm f$, donc un trouve bien
\[
\bm f(\bm x) = \sum_{|\bm \alpha| \leqslant p} \frac{1}{\bm \alpha !} \frac{\partial^{|\bm \alpha|} \bm f}{\partial \bm x^{\bm \alpha}}(\bm x) (\bm x - \bm x_0)^{\bm \alpha} + R_p(\bm x)
\]
cependant le reste de Lagrange est généralement faux puisque ne pas trouver un seul $\theta \in ]0,1[$ tel que
\[
\bm R_p(\bm x) = \sum_{|\bm \alpha| = p + 1} \frac{1}{\bm \alpha !} \frac{\partial^{p+1} \bm f}{\partial \bm x^{\bm \alpha}} (\bm x + \theta (\bm x - \bm x_0))(\bm x - \bm x_0)^{\bm \alpha}
\]
Toutefois, le reste intégrale est quant à lui bien défini,
\[
\bm R_p(\bm x) = \sum_{|\bm \alpha| = p + 1} \frac{(p+1)!}{\bm \alpha!} \int_0^1 (1-s)^p \frac{\partial^{p+1} \bm f}{\partial \bm x^{\bm \alpha}} (\bm x + s (\bm x - \bm x_0))(\bm x - \bm x_0)^{\bm \alpha} ~\mathrm ds
\]